\scanpage{114}
La fonction $p$, considérée comme morphisme
\[
p:E\longrightarrow\mathbb{P}^1_k\tag*{(52)}
\]
précise l'iso
\[
\dot p:E/\{\pm\mathrm{id}_E\}\simeq\text{droite projective type }\mathbb{P}^1_k\tag*{(53)}
\]
par la condition
\[
p(0_E)=\infty\tag*{(54)}
\]
ce qui détermine $p$ modulo transformation affine
\[
p\longmapsto p_1=ap+b\qquad a\in k^*,\ b\in k.\tag*{(55)}
\]
Les points d'ordre $2$
\[
(\unsure{\varepsilon|\xi}_i)_{i\in J},\quad \varepsilon_i\in{}_{2}E(k)^*\tag*{(56)}
\]
sont transformés par $p$ en des points
\[
\zeta_i\in p(\varepsilon_i)\quad\in\mathbb{P}^1_{\deleted{\unsure{k|*}}}(k)\setminus\{\unsure{\infty|0}\}=k\tag*{(57)}
\]
% slanted
\margin{\uline{NB} On obtient, via des $E,p$, \unsure{n'importe} quel triplet de points distincts de $k$}
Ce sont, bien sûr, avec le point $\infty$, les points de ramification de $p$. La normalisation de Weierstrass, possible en car. $\ne3$, \deleted{qui} est
\[
\underbrace{\sum\zeta_i}_{-a_1}=0\tag*{(58)}
\]
elle détermine $p$ modulo transformation
\[
p\longmapsto p_1=ap\qquad a\in k^*.
\]
En général, on pose
\[
\left\{\begin{aligned}
a_1&=-\sum\zeta_i\\
a_2&=\sum\zeta_i\zeta_{\unsure{1|j}}\\
a_3&=\deleted{\unsure{\prod\zeta_i|\zeta_i}}-\prod\zeta_i
\end{aligned}\right.\tag*{(59)}
\]
% archive sheet number: 113
